% chapters/ch_part7_mathematical_formalisation.tex
% Part VII Chapter: Mathematical Formalisation — arXiv preprint, full body
% Source: NRMO_arXiv_Ikeya.pdf (full pdftotext extraction at /tmp/arxiv.txt)

\chapter{Mathematical Formalisation: arXiv Preprint
(Full Reproduction)}
\label{ch:arxiv-formalisation}


\index{arXiv preprint}
\index{Theorem!T1}
\index{Theorem!T2}
\index{Theorem!T3}
\index{Theorem!T4}
\index{Algorithm 1!Phased Parallel}
\nrmobox{Document identification}{
\textbf{Title}: Non-Ruin Maximization Objective (NRMO):
A Formal Framework for Decision-Making Under Absorbing Failure
States.\\[0.3em]
\textbf{Author}: Takashi Ikeya.\\[0.3em]
\textbf{Source}: arXiv preprint --- 2026.\\[0.3em]
\textbf{Status}: Reproduced verbatim from
\texttt{NRMO\_arXiv\_Ikeya.pdf} (uploaded source).
}

% =====================================================================
\section*{Abstract}
\label{sec:arxiv-abstract}

We introduce the Non-Ruin Maximization Objective (NRMO), a
decision-theoretic framework for agents operating in environments
with absorbing failure states. Standard constrained MDPs optimise
expected reward subject to expected constraint satisfaction,
permitting positive ruin probability in individual trajectories.
NRMO places ruin avoidance as a lexicographic priority via a hard
trajectory-level constraint. We prove four results with full rigor:

\begin{itemize}
\item \textbf{(T1)} A ruin dilution theorem showing ruin probability
$q$ bounds the achievable unconditional time-average reward by
$(1 - q) \cdot R_{\max}$.
\item \textbf{(T2)} A CMDP separation theorem via a two-state MDP
with independent per-step ruin probability, correctly computing
finite-horizon ruin probability.
\item \textbf{(T3)} A variational analysis showing the exponential
threshold family achieves sup-norm near-optimality as $\lambda
\to \infty$, with no spurious interior maximiser.
\item \textbf{(T4)} Polynomial-time CVaR tractability for log-concave
distributions.
\end{itemize}

Experiments across three domains demonstrate 57--65\% ruin reduction
versus approximate baselines; limitations of the baseline
implementations are stated explicitly.

\paragraph{Keywords.}
Non-ruin optimisation, absorbing failure states, constrained MDPs,
safe reinforcement learning, ergodicity, CVaR, lexicographic
constraints.

% =====================================================================
\section{Introduction}
\label{sec:arxiv-introduction}

Consider a portfolio manager subject to forced liquidation upon a
30\% drawdown: a textbook absorbing failure state. Standard RL
optimises $\mathbb{E}[\mathrm{return}]$, but once ruin occurs no
future reward accrues. Peters (2019) shows that for multiplicative
growth processes, any policy with positive ruin probability
converges to ruin with probability one under repeated play. This
motivates NRMO's core design: ruin avoidance as a hard constraint,
not a soft penalty.

\paragraph{Contributions.}

\begin{itemize}
\item \textbf{(C1)} Theorem~1: Ruin Dilution --- unconditional
time-average reward is bounded by $(1 - q) \cdot R_{\max}$.
\item \textbf{(C2)} Lemma~1: Constructive feasibility of
$\Pi_{\mathrm{safe}}(T, \varepsilon)$.
\item \textbf{(C3)} Theorem~2: CMDP Separation via explicit two-state
MDP with correct finite-horizon ruin formula.
\item \textbf{(C4)} Theorem~3: Variational near-optimality of
exponential threshold family (corrected; no interior maximiser
exists).
\item \textbf{(C5)} Theorem~4: Polynomial-time CVaR tractability.
\item \textbf{(C6)} Empirical evaluation with explicitly approximated
baselines.
\end{itemize}

% =====================================================================
\section{Related Work}
\label{sec:arxiv-related-work}

Constrained MDPs (Altman, 1999) and Lagrangian methods (Tessler et
al., 2018) guarantee \emph{expected} constraint satisfaction. CPO
(Achiam et al., 2017) provides trust-region updates with constraint
guarantees at each iterate. These methods permit ruin in individual
trajectories; NRMO enforces \emph{trajectory-level} safety.

Risk-sensitive RL (Chow and Pavone, 2014) minimises CVaR of return;
Theorem~4 formalises the connection between CVaR and ruin
probability.

Non-ergodic economics (Peters, 2019; Kelly, 1956) motivates
prioritising time-average over ensemble-average; NRMO generalises
Kelly to arbitrary MDPs.

Robust MDPs (Iyengar, 2005; Nilim and El Ghaoui, 2005) handle
uncertain kernels; NRMO's robust chance constraint uses this
structure with online set estimation.

% =====================================================================
\section{Problem Formulation}
\label{sec:arxiv-formulation}

Let $M = (S, A, P, r, S_F)$ be an MDP with finite state space $S$,
action space $A$, transition kernel $P$, bounded reward
$r: S \times A \to [-R_{\max}, R_{\max}]$, and absorbing failure
states $S_F \subset S$.

\begin{definition}[Absorbing failure state]
\label{def:arxiv-absorbing}
$s_F \in S_F$ satisfies $P(s_F \mid s_F, a) = 1$ and
$r(s_F, a) = 0$ for all $a$. First hitting time:
\begin{equation}
\tau_F := \inf\{t \ge 0 : s_t \in S_F\}.
\end{equation}
\end{definition}

\begin{definition}[Admissible policy set]
\label{def:arxiv-admissible}
\begin{equation}
\Pi_{\mathrm{safe}}(T, \varepsilon) :=
\{\pi : P_\pi(\tau_F \le T) \le \varepsilon\}.
\end{equation}
\end{definition}

\begin{definition}[NRMO objective --- unconditional]
\label{def:arxiv-objective}
\begin{equation}
U(\pi) := \liminf_{T \to \infty}
\frac{1}{T} \mathbb{E}_\pi\!\left[
\sum_{t=0}^{T-1} r(s_t, a_t)
\right].
\end{equation}
\end{definition}

\paragraph{Design choice.}
Defining $U(\pi)$ as the unconditional time-average (not conditioned
on survival) avoids the measure-theoretic complication of
conditioning on a zero-probability event when ruin probability
approaches 1. The ruin dilution theorem (T1) shows that conditioning
is unnecessary: the ruin event itself suppresses $U(\pi)$
multiplicatively. For policies in $\Pi_{\mathrm{safe}}(T, 0)$,
conditional and unconditional averages coincide since
$P(\tau_F = \infty) = 1$.

\paragraph{Robust chance constraint.}
When $P$ is unknown, NRMO uses:

\begin{equation}
\sup_{P \in \mathcal{P}}\,
P_{\pi, P}(\tau_F \le T) \le \varepsilon,
\end{equation}

where $\mathcal{P}$ is a Hoeffding confidence set with radius
$O(n^{-1/2})$.

% =====================================================================
\section{Main Theoretical Results}
\label{sec:arxiv-main-results}

\subsection{Theorem 1: Ruin Dilution}
\label{sec:arxiv-thm-1}

\begin{theorem}[Ruin Dilution]
\label{thm:arxiv-ruin-dilution}
Let $\pi$ be any policy with $P_\pi(\tau_F < \infty) = q \in [0, 1)$
and bounded reward $|r| \le R_{\max}$. Then:
\begin{enumerate}
\item[(i)] $U(\pi) \le (1 - q) \cdot R_{\max}$.
\item[(ii)] Any policy $\pi_0$ with $U(\pi_0) > (1 - q) R_{\max}$
strictly dominates $\pi$: $U(\pi_0) > U(\pi)$.
\end{enumerate}
\end{theorem}

\begin{proof}
\medskip\noindent\textbf{Part (i).}
For any $T$, decompose
$\mathbb{E}_\pi[(1/T) \sum_{t=0}^{T-1} r_t]$ by the partition
$\{\tau_F < \infty\}$ (probability $q$) and $\{\tau_F = \infty\}$
(probability $1 - q$):

\begin{equation}
\mathbb{E}_\pi\!\left[\frac{1}{T} \sum r_t\right]
= q \cdot \mathbb{E}\!\left[
\frac{1}{T} \sum r_t \mid \tau_F < \infty
\right]
+ (1 - q) \cdot \mathbb{E}\!\left[
\frac{1}{T} \sum r_t \mid \tau_F = \infty
\right].
\end{equation}

\medskip\noindent\textbf{Bounding the ruin term.}
On $\{\tau_F < \infty\}$: for all $t \ge \tau_F$, $r(s_t, a_t) = 0$
by Definition~\ref{def:arxiv-absorbing} (absorbing state).
Therefore:

\begin{equation}
\mathbb{E}\!\left[\frac{1}{T} \sum_{t=0}^{T-1} r_t \mid
\tau_F < \infty\right]
= \mathbb{E}\!\left[\frac{1}{T} \sum_{t=0}^{\tau_F - 1} r_t \mid
\tau_F \le T\right]
\le R_{\max} \cdot \mathbb{E}\!\left[\frac{\tau_F}{T} \mid
\tau_F < \infty\right].
\end{equation}

Since $|S|$ is finite and $s_F$ is absorbing,
$\mathbb{E}[\tau_F \mid \tau_F < \infty]$ is finite (bounded hitting
time for irreducible finite chains). Taking $\limsup$ as $T \to
\infty$:

\begin{equation}
\lim_{T \to \infty} R_{\max} \cdot
\mathbb{E}[\tau_F / T \mid \cdot] = 0.
\end{equation}

\medskip\noindent\textbf{Bounding the survival term.}
$\mathbb{E}[(1/T) \sum r_t \mid \tau_F = \infty] \le R_{\max}$ for
all $T$.

\medskip\noindent\textbf{Combining.}
Taking $\limsup$ of both sides:

\begin{equation}
U(\pi) = \liminf_{T \to \infty}
\mathbb{E}_\pi\!\left[\frac{1}{T} \sum r_t\right]
\le \limsup
\le q \cdot 0 + (1 - q) \cdot R_{\max}
= (1 - q) R_{\max}.
\end{equation}

\medskip\noindent\textbf{Part (ii)}
is immediate: if $U(\pi_0) > (1 - q) R_{\max} \ge U(\pi)$, then
$U(\pi_0) > U(\pi)$.
\end{proof}

\begin{remark}[Scope]
\label{rem:arxiv-thm1-scope}
Part~(ii) requires the existence of a policy $\pi_0$ achieving
$U(\pi_0) > (1 - q) R_{\max}$. This is not guaranteed: when safe
actions yield low reward, all zero-ruin policies may satisfy
$U(\pi_0) \le (1 - q) R_{\max}$. The theorem quantifies the
\emph{ceiling} imposed by ruin on any positive-ruin policy; whether
a zero-ruin policy achieves above that ceiling is domain-dependent.
\end{remark}

\begin{remark}[Conditional vs unconditional]
\label{rem:arxiv-thm1-conditional}
The unconditional $U(\pi)$ conflates survival quality with ruin
probability. A policy with $q = 0.5$ and high survival reward may
have $U(\pi) \le 0.5 R_{\max}$. This is the intended behaviour:
NRMO penalises ruin probability in the objective itself, forcing the
agent to jointly optimise survival probability and reward.
\end{remark}

\subsection{Lemma 1: Feasibility}
\label{sec:arxiv-lemma-1}

\begin{lemma}[Feasibility of $\Pi_{\mathrm{safe}}$]
\label{lem:arxiv-feasibility}
Assume:
\begin{itemize}
\item[(A1)] $S$ is finite;
\item[(A2)] there exists a stationary policy $\pi^{\mathrm{hold}}$
satisfying $P_{\pi^{\mathrm{hold}}}(\tau_F \le T) = 0$ for all $T$;
\item[(A3)] $\varepsilon > 0$.
\end{itemize}
Then $\Pi_{\mathrm{safe}}(T, \varepsilon) \ne \emptyset$.

A sufficient condition for (A2): there exists $s^* \in S \setminus
S_F$ and $a^* \in A$ with $P(s^* \mid s^*, a^*) = 1$ (self-loop
action at a safe state).
\end{lemma}

\begin{proof}
Under (A2), $P_{\pi^{\mathrm{hold}}}(\tau_F \le T) = 0 \le
\varepsilon$ for all $T, \varepsilon > 0$. So
$\pi^{\mathrm{hold}} \in \Pi_{\mathrm{safe}}(T, \varepsilon)$. For
the sufficient condition: $\pi^{\mathrm{hold}}(s) = a^*$ for all $s$
keeps all trajectories at $s^* \notin S_F$, satisfying (A2).
\end{proof}

\subsection{Theorem 2: CMDP Separation}
\label{sec:arxiv-thm-2}

\nrmobox{Correction note (from v2)}{
The prior construction
($s_0 \to s_1 \to s_0$ recurrent loop with $s_F$ absorbing from
$s_0$) was flawed: for any $\theta > 0$ in such a recurrent chain,
the long-run ruin probability equals 1 (standard absorbing Markov
chain result). The formula
$q(\theta) = \theta(1 - p) / (1 - \theta p)$ computed only the
probability of ruin on a single visit, not eventually. The
corrected construction uses a two-state MDP with independent
per-step ruin events, for which finite-horizon ruin probability is
exactly computable.
}

\begin{theorem}[NRMO Strictly Separates from Lagrangian CMDP]
\label{thm:arxiv-cmdp-separation}
There exists an explicit MDP $M^*$ and constraint budget
$b \in (0, 1)$ such that:

\begin{enumerate}
\item[(i)] the primal Lagrangian CMDP optimum $\pi^*$ satisfies
$P_{\pi^*}(\tau_F \le T) = b > 0$. Specifically, the unconstrained
optimal policy violates the constraint
$(P(\tau_F \le T) \mid \theta = 1) > b$, so by convexity the
constrained optimum is the unique $\theta^* \in (0, 1)$ satisfying
$\Omega(\theta^*, T) = b$, corresponding to the unique KKT
multiplier
$\lambda^* = (\mathrm{d}\mathbb{E}[\mathrm{return}] /
\mathrm{d}\theta) / (\mathrm{d}\Omega / \mathrm{d}\theta)
\rvert_{\theta = \theta^*} > 0$.

\item[(ii)] $\Pi_{\mathrm{safe}}(T, 0) \ne \emptyset$, i.e., NRMO
achieves zero ruin probability.
\end{enumerate}

Hence on $M^*$, NRMO achieves strictly lower ruin probability than
the CMDP primal optimum.
\end{theorem}

\begin{proof}
\medskip\noindent\textbf{Construction of $M^*$.}
$S = \{s_0, s_F\}$, $A = \{a_{\mathrm{safe}},
a_{\mathrm{risky}}\}$. Parameters: $R = 5$, $p = 0.7$, $T = 10$,
$b = 0.20$. Transitions: from $s_0$, $a_{\mathrm{risky}}$ moves to
$s_F$ with probability $(1 - p) = 0.3$ and stays at $s_0$ with
probability $p = 0.7$, earning $r = R = 5$. From $s_0$,
$a_{\mathrm{safe}}$ stays at $s_0$ with probability 1, $r = 0$.
$s_F$ is absorbing with $r = 0$.

\medskip\noindent\textbf{Key property: independent per-step ruin
events.}
Under stationary mixed policy $\theta = P(a_{\mathrm{risky}} \mid
s_0)$: in any period where $s_t = s_0$, the agent enters $s_F$ with
probability $\theta(1 - p)$ independently of history. Therefore the
finite-horizon ruin probability is:

\begin{equation}
\Omega(\theta, T) := P(\tau_F \le T)
= 1 - (1 - \theta(1 - p))^T.
\end{equation}

This formula is exact: the agent survives step $t$ if and only if it
did not enter $s_F$ at any step $1, \ldots, t$, each event having
probability $1 - \theta(1 - p)$ independently (while at $s_0$, which
is guaranteed given survival so far).

\medskip\noindent\textbf{CMDP analysis.}
The Lagrangian objective is:

\begin{equation}
\mathcal{L}(\theta, \lambda)
= \mathbb{E}[\mathrm{return}] - \lambda(\Omega(\theta, T) - b),
\end{equation}

\begin{equation}
\mathbb{E}[\mathrm{return}] = \theta \cdot R \cdot
\mathbb{E}[\text{steps before ruin}]
= \theta R \cdot T(1 - \theta(1 - p))
\quad (\text{to first order for small ruin probability per step}).
\end{equation}

The CMDP constraint $\Omega(\theta^*, T) = b$ binds at the optimum
for $b \in (0, 1)$. Solving $1 - (1 - \theta^*(1 - p))^T = b$ gives:

\begin{equation}
\theta^* = \frac{1 - (1 - b)^{1/T}}{1 - p}.
\end{equation}

For $T = 10$, $b = 0.20$, $p = 0.7$:
$\theta^* = (1 - (0.80)^{0.1}) / 0.3 = (1 - 0.9777) / 0.3 = 0.0744$.
This gives $P(\tau_F \le 10) = 1 - (1 - 0.0744 \times 0.3)^{10} =
1 - (0.9777)^{10} = 0.200 = b > 0$.

Note that $\lambda$ is not a free parameter we choose; it is the
unique KKT multiplier determined by the primal-dual optimality
conditions. The unconstrained optimum ($\lambda = 0$) is
$\theta = 1$ with $\Omega(1, 10) = 1 - (0.7)^{10} = 0.972 > b =
0.20$, violating the constraint. Since $\mathbb{E}[\mathrm{return}]$
is increasing in $\theta$ and $\Omega$ is increasing in $\theta$
(both strictly), the constraint $\Omega(\theta, T) \le b$ is binding
at the optimum by the complementary slackness condition: the optimal
$\theta^*$ is the unique solution to $\Omega(\theta^*, T) = b$, and
the unique KKT multiplier is:

\begin{equation}
\lambda^* = \frac{\mathrm{d}\mathbb{E}[\mathrm{return}] /
\mathrm{d}\theta}{\mathrm{d}\Omega / \mathrm{d}\theta}
\bigg\rvert_{\theta = \theta^*} > 0.
\end{equation}

Hence the primal CMDP optimum satisfies $P(\tau_F \le T) = b =
0.20 > 0$. For $\lambda > \lambda^*$, the dual solution would push
$\theta$ below $\theta^*$, but this is infeasible since $\theta^*$
is the binding primal solution; for $\lambda < \lambda^*$, the
constraint is violated. The unique primal-dual pair
$(\theta^*, \lambda^*)$ is the CMDP optimum.

\medskip\noindent\textbf{NRMO analysis.}
Setting $\varepsilon = 0$: $\Omega(\theta, 10) = 0$ requires $\theta
= 0$. Lemma~\ref{lem:arxiv-feasibility} applies
($a_{\mathrm{safe}}$ is the hold action). NRMO achieves
$P(\tau_F \le T) = 0$ while CMDP achieves $b = 0.20$.
\end{proof}

\begin{remark}[Interpretation]
\label{rem:arxiv-thm2-interpretation}
The CMDP optimum saturates the ruin budget $b$: it uses every unit
of permitted ruin probability to earn reward. NRMO forfeits all
risky reward to achieve zero ruin. The reward cost is
$U(\pi_{\mathrm{safe}}) = 0$ vs.~$U(\pi^*(\lambda)) =
\theta^* R \times (\text{survival factor}) > 0$. This trade-off is
explicit and domain-dependent.
\end{remark}

\begin{remark}[Generality]
\label{rem:arxiv-thm2-generality}
The separation result holds for all $T \ge 1$, all $b \in (0, 1)$,
and all $R > 0$, $p \in (0, 1)$. The specific numerical example
($T = 10$, $b = 0.20$, $p = 0.7$, $R = 5$) is illustrative; the
construction is valid for the full parameter range.
\end{remark}

\subsection{Theorem 3: Variational Analysis of Exponential
Threshold}
\label{sec:arxiv-thm-3}

\nrmobox{Correction note (from v2)}{
The prior proof applied the Euler-Lagrange equation to
$\Sigma(f) = \int f'(x)(1 - x)\,\mathrm{d}x$, incorrectly yielding a
logarithmic interior maximiser. The correct analysis via integration
by parts shows
$\Sigma(f) = -\varepsilon_{\mathrm{floor}} + \int f(x)\,
\mathrm{d}x$, reducing the problem to maximising $\int f(x)\,
\mathrm{d}x$. This problem has no smooth interior maximiser in
$\mathcal{F}$: the supremum is achieved by a step function at
$x = 0^+$, approached by $f_\lambda$ as $\lambda \to \infty$. The
theorem is restated accordingly.
}

\begin{theorem}[Exponential Family Achieves Variational Supremum]
\label{thm:arxiv-variational}
Let $\mathcal{F} = \{f : [0, 1] \to [\varepsilon_{\mathrm{floor}},
\varepsilon_{\mathrm{global}}]:$ $f$ monotone non-decreasing,
$f(0) = \varepsilon_{\mathrm{floor}}$,
$f(1) = \varepsilon_{\mathrm{global}}$, $f \in C^1\}$. Define the
weighted sensitivity functional
$\Sigma(f) = \int_0^1 f'(x)(1 - x)\,\mathrm{d}x$.

\begin{enumerate}
\item[(i)] Via integration by parts:
$\Sigma(f) = (\varepsilon_{\mathrm{global}} -
\varepsilon_{\mathrm{floor}}) - \int_0^1 f(x)\,\mathrm{d}x + C$,
so maximising $\Sigma(f)$ is equivalent to maximising
$\int f(x)\,\mathrm{d}x$ over $\mathcal{F}$.

\item[(ii)] The infimum of $\int f(x)\,\mathrm{d}x$ over
$\mathcal{F}$ is $\varepsilon_{\mathrm{floor}}$, approached but not
achieved by
$f_\lambda(x) = \varepsilon_{\mathrm{floor}} +
\Delta(1 - e^{-\lambda x}) / (1 - e^{-\lambda})$ as $\lambda \to
\infty$, where $\Delta = \varepsilon_{\mathrm{global}} -
\varepsilon_{\mathrm{floor}}$.

\item[(iii)] The convergence rate is:
$\int f_\lambda(x)\,\mathrm{d}x - \varepsilon_{\mathrm{floor}} =
O(1 / \lambda)$.
\end{enumerate}
\end{theorem}

\begin{proof}
\medskip\noindent\textbf{Part (i): Integration by parts.}
$\Sigma(f) = \int_0^1 f'(x)(1 - x)\,\mathrm{d}x$. Integrate by parts
with $u = (1 - x)$, $\mathrm{d}v = f'(x)\,\mathrm{d}x$:

\begin{equation}
\Sigma(f) = [(1 - x) f(x)]_0^1 - \int_0^1 f(x) \cdot (-1)\,\mathrm{d}x
= (0 \cdot f(1) - 1 \cdot f(0)) + \int_0^1 f(x)\,\mathrm{d}x
= -\varepsilon_{\mathrm{floor}} + \int_0^1 f(x)\,\mathrm{d}x.
\end{equation}

Since $-\varepsilon_{\mathrm{floor}}$ is constant, maximising
$\Sigma(f) \Leftrightarrow$ maximising
$\int f(x)\,\mathrm{d}x$. The maximum of $\int f(x)\,\mathrm{d}x$
over $\mathcal{F}$ is $\sup \int f = \varepsilon_{\mathrm{global}}$
(constant function at upper bound), so we want $f$ as large as
possible. For a monotone function with $f(0) =
\varepsilon_{\mathrm{floor}}$ and $f(1) =
\varepsilon_{\mathrm{global}}$, maximising $\int f(x)\,\mathrm{d}x$
means rising to $\varepsilon_{\mathrm{global}}$ as quickly as
possible from $x = 0$.

\medskip\noindent\textbf{Part (ii): No smooth interior maximiser.}
The supremum of $\int f(x)\,\mathrm{d}x$ over $\mathcal{F}$ is
achieved in the limit by the step function

\begin{equation}
f^*(x) = \varepsilon_{\mathrm{floor}} \cdot \mathbf{1}_{\{x = 0\}}
+ \varepsilon_{\mathrm{global}} \cdot \mathbf{1}_{\{x > 0\}},
\end{equation}

which gives $\int f^* = \varepsilon_{\mathrm{global}}$. But $f^*
\notin C^1(\mathcal{F})$. Therefore no smooth maximiser exists in
$\mathcal{F}$; the supremum is not attained.

The exponential family $f_\lambda$ with large $\lambda$ approximates
$f^*$: $f_\lambda(x) \to \varepsilon_{\mathrm{global}}$
exponentially fast in $\lambda x$ for any $x > 0$. Hence
$\sup_{\lambda > 0} \int f_\lambda = \varepsilon_{\mathrm{global}}
= \sup_{f \in \mathcal{F}} \int f$.

\medskip\noindent\textbf{Part (iii): Convergence rate.}

\begin{equation}
\int_0^1 f_\lambda(x)\,\mathrm{d}x =
\varepsilon_{\mathrm{floor}} + \frac{\Delta}{1 - e^{-\lambda}}
\cdot \int_0^1 (1 - e^{-\lambda x})\,\mathrm{d}x =
\varepsilon_{\mathrm{floor}} +
\frac{\Delta}{1 - e^{-\lambda}}
\cdot \!\left[1 - \frac{1 - e^{-\lambda}}{\lambda}\right].
\end{equation}

\begin{equation}
\int f_\lambda = \varepsilon_{\mathrm{floor}} +
\Delta \cdot \!\left[1 - \frac{1}{\lambda} + O(e^{-\lambda})\right].
\end{equation}

Therefore $\varepsilon_{\mathrm{global}} - \int f_\lambda =
\Delta / \lambda + O(\Delta e^{-\lambda}) = O(\Delta / \lambda)$.
\end{proof}

\begin{remark}[Practical implication]
\label{rem:arxiv-thm3-practical}
The theorem justifies the exponential family as the canonical
approximation to the variational supremum among smooth functions.
Larger $\lambda$ gives better approximation but steeper gradient
near $x = 0$, which can cause numerical instability in discrete
implementations. The choice $\lambda = 3$ is a practical balance,
not a theoretically unique optimum. For $\lambda = 3$:
$\int f_3 - \varepsilon_{\mathrm{floor}} \approx
\varepsilon_{\mathrm{global}} - \Delta / 3$, leaving 33\% of the
sensitivity gain unrealised relative to the limit.
\end{remark}

\subsection{Theorem 4: CVaR Tractability}
\label{sec:arxiv-thm-4}

\begin{theorem}[Polynomial-Time CVaR Bound]
\label{thm:arxiv-cvar}
Let $Z = \mathbf{1}_{\{\tau_F \le T\}} \in \{0, 1\}$ be the ruin
indicator. For any $\alpha \in (0, 1)$:

\begin{enumerate}
\item[(i)] $\mathrm{CVaR}_\alpha(Z) \ge P(Z = 1) =
P(\tau_F \le T)$.

\item[(ii)] The problem
$\min_\pi \mathrm{CVaR}_\alpha(\mathrm{cost}_\pi)$ subject to
$\mathrm{CVaR}_\alpha \le \varepsilon$ is solvable as an LP in time
$O(|S|^2 |A| T \cdot \log(1 / \varepsilon))$ when the loss
distribution is log-concave.
\end{enumerate}

Applicability: geometric stopping times (D1) and log-normal drawdown
(D2) are log-concave.
\end{theorem}

\begin{proof}
\medskip\noindent\textbf{Part (i).}
For $Z \in \{0, 1\}$: $\mathrm{VaR}_\alpha(Z) \in \{0, 1\}$.

\medskip\noindent\textit{Case A: $\alpha \ge 1 - P(Z = 1)$,} so
$\mathrm{VaR}_\alpha(Z) = 1$ and $\mathrm{CVaR}_\alpha(Z) = 1 \ge
P(Z = 1)$.

\medskip\noindent\textit{Case B: $\alpha < 1 - P(Z = 1)$,} so
$\mathrm{VaR}_\alpha(Z) = 0$ and:

\begin{equation}
\mathrm{CVaR}_\alpha(Z) = \frac{1}{1 - \alpha}
\mathbb{E}[Z \cdot \mathbf{1}_{\{Z \ge 0\}}]
= \frac{P(Z = 1)}{1 - \alpha} \ge P(Z = 1).
\end{equation}

In both cases $\mathrm{CVaR}_\alpha(Z) \ge P(Z = 1)$.

\medskip\noindent\textbf{Part (ii).}
By Rockafellar-Uryasev (2000):

\begin{equation}
\mathrm{CVaR}_\alpha(Z) =
\min_v\!\left\{v + \frac{1}{1 - \alpha}
\mathbb{E}[\max(Z - v, 0)]\right\}.
\end{equation}

In the occupancy-measure LP (Altman 1999, Ch.~2),
$\mathbb{E}_\pi[\max(Z - v, 0)]$ is linear in occupancy measures
$\{\mu(s, a, t)\}$ for fixed $v$. The LP has $|S| |A| T + 1$
variables; interior-point methods require $O(|S|^2 |A| T \cdot
\log(1 / \varepsilon))$ operations. Log-concavity ensures the cost
distribution is unimodal, guaranteeing the CVaR bound is tight (the
LP formulation is exact by construction; Altman 1999, Ch.~2).
\end{proof}

% =====================================================================
\section{Algorithm}
\label{sec:arxiv-algorithm}

The Phased Parallel Execution architecture resolves the
contradiction between full parallelism (Phase~1) and Type ZERO
suppression authority (Phase~3 only). The DEADLOCK handler (Step~8)
operationalises Remark~3 of Lemma~1: when all options fail APSOC,
the system recognises near-infeasibility and outputs SAFE\_EXIT
(the hold policy $\pi^{\mathrm{hold}}$).

\paragraph{Algorithm 1: NRMO Execution.}

\begin{lstlisting}[basicstyle=\ttfamily\small,frame=single]
Input:  M, T, eps_global, eps_floor, lambda, B_t
Output: a_t in A or SAFE_EXIT

1: Compute eps_t = eps_floor
                  + (eps_global - eps_floor) * (1 - exp(-lambda*B_t))

2: PHASE 1 [parallel, read-only]:
     module m computes score_m(s_t)

3: PHASE 2 [barrier]:
     await all modules; T_max cutoff;
     late scores set to 0

4: PHASE 3 [serial]:
     Priority Arbiter aggregates ->
     Tri-Option Generator produces A, B, C

5: For o in {A, B, C}:
     verify APSOC (reversibility, MaxLoss, exit trigger)

6: Compute CVaR upper bound on
   P(tau_F <= T) for each o (Theorem 4)

7: O_valid <- {o : CVaR bound <= eps_t and APSOC passed}

8: if O_valid == empty:
     DEADLOCK handler -> return SAFE_EXIT (pi_hold)

9: else:
     return argmax_{o in O_valid} U(pi_o)
\end{lstlisting}

% =====================================================================
\section{Experiments}
\label{sec:arxiv-experiments}

\nrmobox{Correction note --- baseline implementations}{
CPO and RCPO are approximated as analytic greedy Lagrangian policies
($\lambda_{\mathrm{CPO}} = 2.0$, $\lambda_{\mathrm{RCPO}} = 4.0$).
These capture qualitative behaviour but are not faithful
reproductions of the published trust-region or reward-shaping
algorithms. Results should be interpreted as directional evidence.
Comparison against proper library implementations is deferred.
}

\subsection{Domain Descriptions}
\label{sec:arxiv-domains}

\paragraph{D1: Stochastic GridWorld ($5 \times 5$, $n = 500$ runs).}
Absorbing holes at $\{(0, 2), (1, 4), (2, 1), (3, 3), (4, 1)\}$.
Noisy hole-proximity observation (policy-specific noise
$\sigma = 0.04 / 0.12 / 0.09 / 0.16$). Slip probability 5\%.
Reward: $+10$ at goal, $-5$ at hole, $-0.02$/step.

\paragraph{D2: Portfolio Management (3-year, $n = 300$ runs).}
Risky asset $\mathcal{N}(0.035\%, 1.0\%)$/day with 3\%
crash-regime probability ($\mathcal{N}(-2.5\%, 2.0\%)$). Ruin: 25\%
peak drawdown. NRMO allocation: $20\% + 25\% \times B_t$;
baselines: fixed 45--60\%.

\paragraph{D3: Inventory Control (200 periods, $n = 300$ runs).}
Demand $\mathcal{N}(20, 8)$. Ruin: 5 consecutive stockout periods.
NRMO triggers emergency restock at streak~$\ge 2$.

\subsection{Results}
\label{sec:arxiv-results}

\begin{table}[ht]
\centering
\scriptsize
\begin{tabularx}{\linewidth}{lXXXXXXXX}
\toprule
\textbf{Method} & \textbf{D1 Reward} & \textbf{D1 Ruin$\downarrow$}
& \textbf{D2 Return} & \textbf{D2 Ruin$\downarrow$} &
\textbf{D2 Sharpe} & \textbf{D3 Profit} &
\textbf{D3 Ruin$\downarrow$} & \textbf{D3 Service} \\
\midrule
\textbf{NRMO} & \textbf{8.22} & \textbf{0.106} & $-9.0\%$
& \textbf{0.193} & $-0.49$ & \textbf{17{,}548} & 0.000 & 0.997 \\
CPO & 6.13 & 0.248 & $-11.3\%$ & 0.467 & $-0.60$ & 16{,}306 & 0.000
& 0.998 \\
RCPO & 8.06 & 0.118 & $-9.7\%$ & 0.250 & $-0.54$ & 17{,}176 & 0.000
& 0.995 \\
Lagrangian & 5.69 & 0.278 & $-13.8\%$ & 0.550 & $-0.77$ & 16{,}469
& 0.000 & 0.976 \\
\bottomrule
\end{tabularx}
\caption[arXiv experimental results]{Experimental results.
$\downarrow$ = lower is better. Baselines are analytic
approximations (see caveat above).}
\label{tab:arxiv-empirical}
\end{table}

\paragraph{D1.}
NRMO reduces ruin by 57\% vs.~Lagrangian ($0.106$ vs.~$0.278$) and
10\% vs.~RCPO ($0.118$). The hard ruin filter rejects options with
CVaR estimate above $\varepsilon_t = 0.10$ entirely, while
Lagrangian penalises ruin softly.

\paragraph{D2.}
Buffer-adaptive allocation reduces NRMO ruin to $0.193$ vs.~$0.550$
for Lagrangian ($-65\%$). As drawdown approaches the threshold,
$\varepsilon_t$ tightens ($B_t \to 0 \Rightarrow \varepsilon_t \to
\varepsilon_{\mathrm{floor}}$) and risky allocation decreases
mechanically.

\paragraph{D3.}
All methods achieve zero ruin; NRMO leads on profit ($+7.6\%$
vs.~CPO) due to the emergency restock trigger preventing stockout
cascades.

% =====================================================================
\section{Discussion}
\label{sec:arxiv-discussion}

\paragraph{T1 scope.}
Theorem~1 bounds the achievable unconditional time-average reward of
any policy with ruin probability $q$. It does not assert that NRMO
achieves higher reward than CMDP in general: when safe actions earn
zero reward (as in $M^*$ of Theorem~2), NRMO's zero-ruin policy has
$U = 0$ while CMDP achieves $U > 0$. The practitioner's choice to
use NRMO is the assertion that zero ruin is worth this cost.

\paragraph{T2 scope.}
The separation instance $M^*$ is simple by design to enable clean
analysis. In environments where CMDP's ruin probability is
negligible ($\theta^*(1 - p)^T \ll 1$), the practical difference
between NRMO and CMDP vanishes. The separation is theoretically
strict but practically relevant only when $T \cdot \theta^*(1 - p)
= O(1)$.

\paragraph{T3 scope.}
The variational result identifies the exponential family as the
canonical smooth approximation to the (unattainable) step-function
supremum. It does not provide a unique optimal $\lambda$: larger
$\lambda$ is always better from the variational perspective,
constrained only by implementation numerics.

\paragraph{Open problems.}

\begin{enumerate}
\item Extension to continuous state spaces via function
approximation.
\item POMDP adaptation of the APSOC gate.
\item Online estimation of ambiguity set $\mathcal{P}$ with
finite-sample guarantees.
\item Comparison against proper CPO/RCPO implementations using
Safety Gym benchmarks.
\end{enumerate}

% =====================================================================
\section{Conclusion}
\label{sec:arxiv-conclusion}

We presented NRMO with four theorems proved with full mathematical
rigour, correcting three errors identified in the prior draft.
Theorem~1 provides a clean multiplicative bound on achievable reward
as a function of ruin probability, using unconditional time-average
to avoid measure-theoretic subtlety. Theorem~2 proves CMDP
separation via an explicit two-state MDP with exact finite-horizon
ruin formula $\Omega(\theta, T) = 1 - (1 - \theta(1 - p))^T$.
Theorem~3 correctly identifies that no smooth interior maximiser
exists for the sensitivity functional, and that the exponential
family approaches the variational supremum at rate $O(1 / \lambda)$.
Theorem~4 establishes polynomial-time CVaR tractability. Experiments
with explicitly approximated baselines provide directional empirical
support.

% =====================================================================
\section*{References}
\label{sec:arxiv-references}

\begin{enumerate}
\item Achiam, J., Held, D., Tamar, A., \& Abbeel, P.~(2017).
Constrained Policy Optimization. \textit{ICML 2017}.

\item Altman, E.~(1999). \textit{Constrained Markov Decision
Processes}. Chapman \& Hall/CRC.

\item Blanchini, F.~(1999). Set invariance in control.
\textit{Automatica}, 35(11), 1747--1767.

\item Chow, Y., \& Pavone, M.~(2014). A unified approach to
risk-sensitive MDPs. \textit{IEEE CDC 2014}.

\item Chow, Y., Nachum, O., Duez-Guarneri, E., \& Ghavamzadeh,
M.~(2018). A Lyapunov-based approach to safe RL.
\textit{NeurIPS 2018}.

\item Garcia, J., \& Fernandez, F.~(2015). A comprehensive survey on
safe RL. \textit{JMLR}, 16(1), 1437--1480.

\item Iyengar, G.~N.~(2005). Robust dynamic programming.
\textit{Math.\ Operations Research}, 30(2), 257--280.

\item Kelly, J.~L.~(1956). A new interpretation of information rate.
\textit{Bell System Technical Journal}, 35(4), 917--926.

\item Nilim, A., \& El~Ghaoui, L.~(2005). Robust control of Markov
decision processes. \textit{Operations Research}, 53(5),
780--798.

\item Peters, O.~(2019). The ergodicity problem in economics.
\textit{Nature Physics}, 15(12), 1216--1221.

\item Rockafellar, R.~T., \& Uryasev, S.~(2000). Optimization of
conditional value-at-risk. \textit{Journal of Risk}, 2(3),
21--41.

\item Tessler, C., Mankowitz, D.~J., \& Mannor, S.~(2018). Reward
constrained policy optimization. \textit{ICLR 2019}.
\end{enumerate}

% =====================================================================
\section*{Connections to Other Parts}

\begin{itemize}
\item This chapter is the arXiv-format paper. The same theorems are
also presented in monograph form in
Chapter~\ref{ch:theoretical-foundation-v4}.
\item Theorem~\ref{thm:arxiv-ruin-dilution} (Ruin Dilution) is the
mathematical foundation for the v2.0 invariant
\textit{Non-Ruin Constraint}
(Section~\ref{sec:invariants-unchanged}).
\item Lemma~\ref{lem:arxiv-feasibility} (Feasibility) provides the
existence guarantee for SAFE\_EXIT in the engineering implementation
(Part~X).
\item Algorithm~1 (Phased Parallel Execution) is implemented as the
governance pipeline in $\Omega$ Full v5.0 (Part~IX).
\item The empirical results
(Section~\ref{sec:arxiv-results}) at the algorithmic level
complement the civilisation-scale empirical validation in Part~XI.
\end{itemize}
